\documentclass[12pt,a4paper,oneside]{book}
% \usepackage[latin1]{inputenc}
\usepackage{amsmath}
\usepackage{amsfonts}
\usepackage{amssymb}
\usepackage{graphicx}
\usepackage[left=2.5cm,right=2.5cm,top=3cm,bottom=2.5cm]{geometry}
% استفاده از بسته‌ی زیر الزامی نیست ولی با استفاده از آن می‌توانید لینکهای رنگی به مراجع خود داشته باشید. 
\usepackage[colorlinks,citecolor=blue]{hyperref}
\usepackage{xepersian}
\author{پردیس دانشکده‌های فنی}
\title{نمونه یک پایان نامه در لاتک}

\settextfont[Scale=1]{XB Zar}
\setdigitfont[Scale=1]{XB Zar}

%----------------------------------------------------------------------
% شماره‌گذاری به سبک پایان‌نامه‌های فارسی: ۱-۱ به جای ۱.۱
%----------------------------------------------------------------------
\renewcommand{\thesection}{\thechapter-\arabic{section}}
\renewcommand{\thesubsection}{\thesection-\arabic{subsection}}
\renewcommand{\theequation}{\thechapter-\arabic{equation}}
\renewcommand{\thetable}{\thechapter-\arabic{table}}
\renewcommand{\thefigure}{\thechapter-\arabic{figure}}

% اجازه کشش بیشتر به فاصله‌ها برای کاهش سرریز سطرها در متن فارسی
\setlength{\emergencystretch}{3em}